Let $A\subseteq\mathbb{N}$. For an integer $n\ge 2$, define the number of \emph{unordered} representations of $n$ as a sum of two elements of $A$ by
\[
 r_A(n):=|\{\{a,b\}: a\in A,\ b\in A,\ a\le b,\ a+b=n\}|.
\]
(So $a=b$ is allowed; it contributes one representation.)
Define
\[
 B:=\{n\in\mathbb{N}: r_A(n)=1\}.
\]
The questions are:
\begin{enumerate}
\item[(Q1)] Is it true that for every $\varepsilon>0$ there exists $N_0(\varepsilon)$ such that for all $N\ge N_0(\varepsilon)$,
\[
|\{1,\dots,N\}\setminus B| \gg_{\varepsilon} N^{1/2-\varepsilon}?
\]
\item[(Q2)] Is it possible that $|\{1,\dots,N\}\setminus B|=o(N^{1/2})$ as $N\to\infty$ for some choice of $A$?
\end{enumerate}

Ambiguity note: some authors count ordered representations $(a,b)$ instead of unordered, or forbid $a=b$. The qualitative questions change only by constant factors in typical regimes, but for small $n$ the distinction matters; I proceed with the (common) convention above.

2)
